% GENERATED FILE. Edit canonical lessons, then run npm run generate:books.
% canonical-source-hash: fnv1a64:f5dcbe945a76e31a
% canonical-lessons: SA-C456-ksubhyati, SA-C456-prasadayati, SA-C456-samahitah, SA-C456-mardayati, SA-C456-ekantah

\chapter{Calm and Composure}
\label{ch:sa-calm-and-composure}
\hlchaptermodality{456}

\begin{chapteropening}
\textbf{By the end of this chapter:} \emph{I can say how someone stays calm.}

Complete the last lesson of chapter 456: I can say how someone stays calm.
\end{chapteropening}

\section[\texorpdfstring{\sk{क्षुभ्यति} (kṣubhyati)}{क्षुभ्यति (kṣubhyati)}]{\sk{क्षुभ्यति} (kṣubhyati) — is shaken}

\label{lesson:SA-C456-ksubhyati}



\begin{quote}
Before the new one: say the Sanskrit for attachment, then the Sanskrit for tenderness for a child.
\end{quote}

\subsection*{You'll want to know: \sk{क्षुभ्यति}}
\textbf{\sk{क्षुभ्यति}} — \emph{kṣubhyati} — \textquotedblleft{}is shaken\textquotedblright{}.

\begin{quote}
\sk{वार्तया} \sk{जनाः} \sk{क्षुभ्यन्ति।}
\end{quote}

\begin{quote}
\emph{vārtayā janāḥ kṣubhyanti.}
\end{quote}

\begin{quote}
People are shaken by the news.
\end{quote}

Past (he …): \sk{सः} \sk{वार्तया} \sk{क्षुभितवान्}, \emph{saḥ vārtayā kṣubhitavān}.

\begin{cousinweb}
\textbf{\sk{क्षुभ्यति}} (\emph{kṣubhyati}) comes from the root √kṣubh, \textquotedblleft{}to shake, to be agitated\textquotedblright{}; MW's ppp \emph{kṣubhita} is \textquotedblleft{}agitated, disturbed, frightened\textquotedblright{} (MW).
\end{cousinweb}

\begin{grammarlens}[title={What you've built}]
One more word for describing what you lived and felt.
\end{grammarlens}

\subsection*{Your turn}
\begin{itemize}
\raggedright
  \item \emph{Say it:} \emph{kṣubhyati}
  \item \emph{Say it:} \emph{kṣubhyati}, once more
  \item \emph{Say it:} \emph{vārtayā janāḥ kṣubhyanti.} — People are shaken by the news.
  \item \emph{Recall:} say the Sanskrit for attachment, then the Sanskrit for tenderness for a child, then say \emph{kṣubhyati} again
\end{itemize}

\begin{tcolorbox}[breakable,colback=teal!4,colframe=teal!35!black,title={Before you move on}]
What does \sk{क्षुभ्यति} mean? (\textquotedblleft{}Is shaken\textquotedblright{}.) Say it once more.
\end{tcolorbox}
\section[\texorpdfstring{\sk{प्रसादयति} (prasādayati)}{प्रसादयति (prasādayati)}]{\sk{प्रसादयति} (prasādayati) — appeases}

\label{lesson:SA-C456-prasadayati}



\begin{quote}
Before the new one: say the Sanskrit for tenderness for a child, then the Sanskrit for is shaken.
\end{quote}

\subsection*{You'll want to know: \sk{प्रसादयति}}
\textbf{\sk{प्रसादयति}} — \emph{prasādayati} — \textquotedblleft{}appeases\textquotedblright{}.

\begin{quote}
\sk{पुत्रः} \sk{कुपितं} \sk{पितरं} \sk{प्रसादयति।}
\end{quote}

\begin{quote}
\emph{putraḥ kupitaṃ pitaraṃ prasādayati.}
\end{quote}

\begin{quote}
The son calms his angry father.
\end{quote}

Past (he …): \sk{सः} \sk{कुपितं} \sk{पितरं} \sk{प्रसादितवान्}, \emph{saḥ kupitaṃ pitaraṃ prasāditavān}.

\begin{cousinweb}
\textbf{\sk{प्रसादयति}} (\emph{prasādayati}) is built on \emph{pra} and the root √sad, \textquotedblleft{}to sit\textquotedblright{}: \emph{pra-√sad} is \textquotedblleft{}to become placid, to be pleased\textquotedblright{}, and this causative is \textquotedblleft{}to make calm\textquotedblright{}; MW's ppp \emph{prasādita} is \textquotedblleft{}pleased, conciliated\textquotedblright{} (MW).
\end{cousinweb}

\begin{grammarlens}[title={What you've built}]
One more word for describing what you lived and felt.
\end{grammarlens}

\subsection*{Your turn}
\begin{itemize}
\raggedright
  \item \emph{Say it:} \emph{prasādayati}
  \item \emph{Say it:} \emph{prasādayati}, once more
  \item \emph{Say it:} \emph{putraḥ kupitaṃ pitaraṃ prasādayati.} — The son calms his angry father.
  \item \emph{Recall:} say the Sanskrit for tenderness for a child, then the Sanskrit for is shaken, then say \emph{prasādayati} again
\end{itemize}

\begin{tcolorbox}[breakable,colback=teal!4,colframe=teal!35!black,title={Before you move on}]
What does \sk{प्रसादयति} mean? (\textquotedblleft{}Appeases\textquotedblright{}.) Say it once more.
\end{tcolorbox}
\section[\texorpdfstring{\sk{समाहितः} (samāhitaḥ)}{समाहितः (samāhitaḥ)}]{\sk{समाहितः} (samāhitaḥ) — composed; collected}

\label{lesson:SA-C456-samahitah}



\begin{quote}
Before the new one: say the Sanskrit for is shaken, then the Sanskrit for appeases.
\end{quote}

\subsection*{You'll want to know: \sk{समाहितः}}
\textbf{\sk{समाहितः}} — \emph{samāhitaḥ} — \textquotedblleft{}composed; collected\textquotedblright{}.

\begin{quote}
\sk{परीक्षायाम्} \sk{अपि} \sk{छात्रः} \sk{समाहितः} \sk{अस्ति।}
\end{quote}

\begin{quote}
\emph{parīkṣāyām api chātraḥ samāhitaḥ asti.}
\end{quote}

\begin{quote}
Even in the exam the student stays composed.
\end{quote}

\begin{cousinweb}
\textbf{\sk{समाहितः}} (\emph{samāhitaḥ}) is built from \emph{sam}, \emph{ā} and \emph{hita}, \textquotedblleft{}put\textquotedblright{}, from the root √dhā, \textquotedblleft{}to put\textquotedblright{}: put together; MW gives \textquotedblleft{}put or held together\textquotedblright{} and then \textquotedblleft{}composed, collected, concentrated\textquotedblright{} (MW).
\end{cousinweb}

\begin{grammarlens}[title={What you've built}]
One more word for describing what you lived and felt.
\end{grammarlens}

\subsection*{Your turn}
\begin{itemize}
\raggedright
  \item \emph{Say it:} \emph{samāhitaḥ}
  \item \emph{Say it:} \emph{samāhitaḥ}, once more
  \item \emph{Say it:} \emph{parīkṣāyām api chātraḥ samāhitaḥ asti.} — Even in the exam the student stays composed.
  \item \emph{Recall:} say the Sanskrit for is shaken, then the Sanskrit for appeases, then say \emph{samāhitaḥ} again
\end{itemize}

\begin{tcolorbox}[breakable,colback=teal!4,colframe=teal!35!black,title={Before you move on}]
What does \sk{समाहितः} mean? (\textquotedblleft{}Composed; collected\textquotedblright{}.) Say it once more.
\end{tcolorbox}
\section[\texorpdfstring{\sk{मर्दयति} (mardayati)}{मर्दयति (mardayati)}]{\sk{मर्दयति} (mardayati) — rubs; massages}

\label{lesson:SA-C456-mardayati}



\begin{quote}
Before the new one: say the Sanskrit for appeases, then the Sanskrit for composed.
\end{quote}

\subsection*{You'll want to know: \sk{मर्दयति}}
\textbf{\sk{मर्दयति}} — \emph{mardayati} — \textquotedblleft{}rubs; massages\textquotedblright{}.

\begin{quote}
\sk{माता} \sk{बालस्य} \sk{पादौ} \sk{मर्दयति।}
\end{quote}

\begin{quote}
\emph{mātā bālasya pādau mardayati.}
\end{quote}

\begin{quote}
Mother rubs the child's feet.
\end{quote}

Past (she …): \sk{सा} \sk{बालस्य} \sk{पादौ} \sk{मर्दितवती}, \emph{sā bālasya pādau marditavatī}.

\begin{cousinweb}
\textbf{\sk{मर्दयति}} (\emph{mardayati}) is built on the causative of the root √mṛd, \textquotedblleft{}to press, to crush\textquotedblright{}; MW gives \textquotedblleft{}to press or squeeze hard\textquotedblright{} and \textquotedblleft{}to rub\textquotedblright{} (MW).
\end{cousinweb}

\begin{grammarlens}[title={What you've built}]
One more word for describing what you lived and felt.
\end{grammarlens}

\subsection*{Your turn}
\begin{itemize}
\raggedright
  \item \emph{Say it:} \emph{mardayati}
  \item \emph{Say it:} \emph{mardayati}, once more
  \item \emph{Say it:} \emph{mātā bālasya pādau mardayati.} — Mother rubs the child's feet.
  \item \emph{Recall:} say the Sanskrit for appeases, then the Sanskrit for composed, then say \emph{mardayati} again
\end{itemize}

\begin{tcolorbox}[breakable,colback=teal!4,colframe=teal!35!black,title={Before you move on}]
What does \sk{मर्दयति} mean? (\textquotedblleft{}Rubs; massages\textquotedblright{}.) Say it once more.
\end{tcolorbox}
\section[\texorpdfstring{\sk{एकान्तः} (ekāntaḥ)}{एकान्तः (ekāntaḥ)}]{\sk{एकान्तः} (ekāntaḥ) — a secluded place; solitude}

\label{lesson:SA-C456-ekantah}



\begin{quote}
Before the new one: say the Sanskrit for composed, then the Sanskrit for massages.
\end{quote}

\subsection*{You'll want to know: \sk{एकान्तः}}
\textbf{\sk{एकान्तः}} — \emph{ekāntaḥ} — \textquotedblleft{}a secluded place; solitude\textquotedblright{}.

\begin{quote}
\sk{एकान्ते} \sk{मनः} \sk{शान्तं} \sk{भवति।}
\end{quote}

\begin{quote}
\emph{ekānte manaḥ śāntaṃ bhavati.}
\end{quote}

\begin{quote}
In a quiet place the mind grows calm.
\end{quote}

A masculine noun.

\begin{cousinweb}
\textbf{\sk{एकान्तः}} (\emph{ekāntaḥ}) is built from \emph{eka}, \textquotedblleft{}one\textquotedblright{}, and \emph{anta}, \textquotedblleft{}an end, a side\textquotedblright{}: a place set to one side; MW gives \textquotedblleft{}a lonely or retired or secret place\textquotedblright{} (MW).
\end{cousinweb}

\begin{grammarlens}[title={What you've built}]
One more word for describing what you lived and felt.
\end{grammarlens}

\subsection*{Your turn}
\begin{itemize}
\raggedright
  \item \emph{Say it:} \emph{ekāntaḥ}
  \item \emph{Say it:} \emph{ekāntaḥ}, once more
  \item \emph{Say it:} \emph{ekānte manaḥ śāntaṃ bhavati.} — In a quiet place the mind grows calm.
  \item \emph{Recall:} say the Sanskrit for composed, then the Sanskrit for massages, then say \emph{ekāntaḥ} again
\end{itemize}

\begin{tcolorbox}[breakable,colback=teal!4,colframe=teal!35!black,title={Before you move on}]
What does \sk{एकान्तः} mean? (\textquotedblleft{}A secluded place; solitude\textquotedblright{}.) Say it once more.
\end{tcolorbox}
